\section{Resource-normalized measurement and the bi-factor law}
\label{sec:model}

\subsection{The \RNM{} protocol}

If $\n$ validator processes share a host offering $Q$ cores under a fair
scheduler, each receives $Q/\n$ cores, so the per-validator compute budget
shrinks as $\n$ grows. Measured throughput then falls for two unrelated
reasons --- growing consensus communication and a shrinking compute budget ---
and an exponent fitted to such data is a property of the host, not of the
protocol.

\begin{Def}[Resource-normalized measurement]
\label{def:rnm}
A campaign satisfies the \emph{\RNM{} protocol} with quota~$\q$ if every
validator container is granted the same CPU quota~$\q$ and memory limit
independently of~$\n$, and the aggregate reservation $\n\q$ stays below host
capacity with headroom. The \emph{normalized throughput} is
$\TPSn(\n,\q)\eqdef\TPS(\n,\q)/\q$.
\end{Def}

This is implemented with control-group quotas in a container runtime and
requires no modification of the consensus software.

\begin{Assumption}[Separable compute budget]
\label{as:separable}
On the operating range considered,
$\TPS(\n,\cc,\q)=\q\cdot\Phi(\n,\cc)\cdot(1+O(\q))$: the per-validator compute
budget enters multiplicatively and does not interact with the validator count to
first order.
\end{Assumption}

Everything downstream rests on this. It should hold because a validator's work
per block --- message verification and transaction execution --- is CPU-bound
and throttled by the same cgroup quota, which delivers $\q$ core-seconds per
wall second regardless of how many siblings exist: the validator count sets
\emph{how much} work there is, the quota \emph{how fast} it is retired. It
should fail when fixed per-process overhead such as garbage collection consumes
a share of $\q$ that does not shrink with $\q$ (the residual $O(\q)$ term, which
dominates below some quota floor), when $\n\q$ approaches host capacity and
contention reappears through memory bandwidth, cache or the network stack, or
when block production is paced by a fixed period rather than by compute. It is a
working hypothesis with a stated domain of validity, not a law.

\begin{State}[$\q$-invariance of the consensus exponent]
\label{st:qinvariance}
Under Assumption~\ref{as:separable} and the model
$\Phi(\n,\cc)=\Tzero\,\gfun(\cc)\,\n^{-\beta}$, the least-squares estimate of
$\beta$ obtained by regressing $\log\TPSn$ on $\log\n$ at fixed~$\cc$ does not
depend on the quota~$\q$ used in the campaign.
\end{State}

\begin{Proof}
$\log\TPSn(\n,\q)=\log\Phi(\n,\cc)+O(\q)$ by Definition~\ref{def:rnm} and
Assumption~\ref{as:separable}. At fixed~$\q$ the term $O(\q)$ is constant across
the regression in $\log\n$ and is absorbed by the intercept, so the slope equals
$-\beta$ for every~$\q$.
\end{Proof}

The quota thus becomes an experimental control: a small~$\q$ lowers absolute
throughput but leaves the quantity of interest unchanged, which is what allows a
validator set far larger than the number of physical cores to be studied on one
machine. Two falsifiable tests follow. Fitting $\log\TPSn$ on $\log\n$ across
three quota levels with quota dummies and quota\,$\times\log\n$ interactions,
invariance is the hypothesis that the interaction coefficients vanish, tested by
a partial $F$-test; and raw throughput at fixed $(\n,\cc)$ should be
proportional to $\q$ with zero intercept, a positive intercept indicating a
latency-bound regime and a negative one fixed overhead. If either test rejects,
the response is to restrict the protocol to the quota interval on which
invariance survives and to report the exponent together with that interval.

\subsection{The bi-factor law and its identifiability}

Two mechanisms govern throughput: more independent submitters~$\cc$ fill the
pool and the block, raising throughput until saturation, whereas more
validators~$\n$ add communication rounds and signature verification, lowering
it. We model them multiplicatively.

\begin{Def}[Bi-factor scaling law]
\label{def:bifactor}
For a permissioned BFT network under the \RNM{} protocol,
\begin{equation}
\TPS(\n,\cc)=\Tzero\cdot\gfun(\cc)\cdot\n^{-\beta},
\qquad
\gfun(\cc)=\frac{\cc^{\gamma}}{1+\bigl(\cc/\csat\bigr)^{\theta}} ,
\label{eq:bifactor}
\end{equation}
with $\Tzero>0$, $\beta>0$, $\gamma\in(0,1]$, $\theta>\gamma$, $\csat>0$.
\end{Def}

The constraint $\beta>0$ encodes replicated-state-machine semantics: every
validator executes every transaction, so adding validators cannot raise
aggregate throughput. The constraint $\theta>\gamma$ makes $\gfun$ unimodal,
reproducing the collapse once the pool overflows. In the unsaturated regime
$\cc\ll\csat$ the law is log-linear in both factors,
\begin{equation}
\log\TPS=\log\Tzero+\gamma\log\cc-\beta\log\n ,
\label{eq:loglinear}
\end{equation}
and it contains the coherency-dominated regime of the Universal Scalability
Law as the case $\beta=1$, while allowing batching and pipelining to give
$\beta<1$.

Single-factor studies report $\TPS(\n)\propto\n^{\alpha}$. How such a fit can
yield $\alpha>0$ despite $\beta>0$ is simple to state: an experimenter doubles
the network and, quite reasonably, doubles the load generator with it, so
throughput moves down because consensus costs more and up because more
transactions are offered. A regression on $\log\n$ alone cannot separate the two
and attributes the net movement to $\n$. The theorem supplies the exact
bookkeeping.

\begin{Theorem}[Identifiability of the consensus exponent]
\label{thm:identifiability}
Let $\{(\n_{i},\cc_{i},\TPS_{i})\}_{i=1}^{m}$ be generated
by~\eqref{eq:loglinear} with additive noise $\epsilon_{i}$,
$\E[\epsilon_{i}\mid\n_{i}]=0$, $\Var(\epsilon_{i})=\sigma^{2}<\infty$, and let
the campaign follow a \emph{concurrency path}
\begin{equation}
\log\cc_{i}=\lambda\log\n_{i}+\mu+\eta_{i},
\qquad
\E[\eta_{i}\mid\n_{i}]=0,\quad \Var(\eta_{i})<\infty ,
\label{eq:path}
\end{equation}
for some real $\lambda$. Assume that, with $x_{i}=\log\n_{i}$, the empirical
moments converge: $\bar{x}_{m}\to\bar{x}$ and
$m^{-1}\sum_{i}(x_{i}-\bar{x}_{m})^{2}\to s_{x}^{2}>0$, which holds for any
design revisiting a fixed finite set of at least two validator counts with
stable proportions. Let $\ahat_{m}$ be the least-squares slope from regressing
$\log\TPS_{i}$ on $x_{i}$ alone. Then
\begin{equation}
\ahat_{m}\ \xrightarrow{\ \Prob\ }\ \gamma\lambda-\beta .
\label{eq:confounding}
\end{equation}
In particular $\ahat>0$ if and only if $\lambda>\beta/\gamma$, and
$\ahat=-\beta$ if and only if $\lambda=0$, \ie{} only when client concurrency is
held fixed.
\end{Theorem}

\begin{Proof}
Substituting~\eqref{eq:path} into~\eqref{eq:loglinear} and writing
$y_{i}=\log\TPS_{i}$,
\begin{equation}
y_{i}=\underbrace{\bigl(\log\Tzero+\gamma\mu\bigr)}_{\eqdef\,a}
+\bigl(\gamma\lambda-\beta\bigr)x_{i}
+\underbrace{\gamma\eta_{i}+\epsilon_{i}}_{\eqdef\,\xi_{i}} ,
\label{eq:reduced}
\end{equation}
a correctly specified simple linear regression with slope $\gamma\lambda-\beta$,
$\E[\xi_{i}\mid x_{i}]=0$ and
$\Var(\xi_{i})\le2\gamma^{2}\Var(\eta_{i})+2\sigma^{2}<\infty$. Inserting
\eqref{eq:reduced} into the least-squares slope, the constant $a$ cancels
against the centring and $\ahat_{m}=(\gamma\lambda-\beta)+N_{m}/D_{m}$ with
$N_{m}=m^{-1}\sum_{i}(x_{i}-\bar{x}_{m})(\xi_{i}-\bar{\xi}_{m})$ and
$D_{m}=m^{-1}\sum_{i}(x_{i}-\bar{x}_{m})^{2}$. The summands of $N_{m}$ have zero
mean by the law of iterated expectations and finite variance, so
$N_{m}\to0$ in probability while $D_{m}\to s_{x}^{2}>0$; the continuous mapping
theorem applied to the ratio yields~\eqref{eq:confounding}. Since $\gamma>0$,
$\gamma\lambda-\beta>0\iff\lambda>\beta/\gamma$ and
$\gamma\lambda-\beta=-\beta\iff\lambda=0$.
\end{Proof}

\begin{Remark}[Order of magnitude]
The bias is not a second-order correction. With $\gamma=0.8$, $\beta=0.5$ and a
load generator scaled proportionally to the network ($\lambda=1$), the theorem
gives $\ahat=+0.3$: the same measurements that contain a genuine $-0.5$ decay
report growth. Extrapolated from $\n=5$ to $\n=60$ the two exponents differ by
a factor $12^{0.8}\approx7.3$, which is the order of the discrepancy we set out
to explain, and which we conjecture is the mechanism behind previously reported
positive exponents including our own~\cite{peniak2026a}.
\end{Remark}

\begin{Corollary}[Non-transferability]
\label{cor:sign}
Two laboratories measuring the same implementation along paths
$\lambda_{1}\ne\lambda_{2}$ obtain exponents differing by
$\gamma(\lambda_{1}-\lambda_{2})$, and predictions extrapolated to $\n_{t}$
differ by the factor $\n_{t}^{\gamma(\lambda_{1}-\lambda_{2})}$. Published
exponents are therefore not comparable unless the concurrency paths are
reported.
\end{Corollary}

Estimation is straightforward. On a factorial grid the unsaturated subset is
fitted by ordinary least squares in log-space,
\begin{equation}
\min_{\log\Tzero,\gamma,\beta}\ \sum_{i=1}^{m}
\bigl(\log\TPS_{i}-\log\Tzero-\gamma\log\cc_{i}+\beta\log\n_{i}\bigr)^{2} ,
\label{eq:ols}
\end{equation}
a two-regressor problem with a closed-form solution; $(\csat,\theta)$ then
follow from nonlinear least squares on the full $\cc$-sweep, and $\hat\sigma$ is
the standard deviation of the log-space residuals. A factorial design guarantees
$\Cov(\log\n,\log\cc)=0$, so the effects are orthogonal by construction and
Theorem~\ref{thm:identifiability} applies with $\lambda=0$.
